\subsection{Squared-error benefit of posterior-mean propagation}

Let \(T\) be a latent temperature with finite second moment and let \(Y\)
denote the available observation information. For any estimator \(a(Y)\),
conditioning on \(Y\) gives the orthogonal decomposition
\[
\mathbb{E}\!\left[(T-a(Y))^2\mid Y\right]
=
\operatorname{Var}(T\mid Y)
+
\left(a(Y)-\mathbb{E}[T\mid Y]\right)^2.
\]
Therefore \(a^*(Y)=\mathbb{E}[T\mid Y]\) minimizes conditional squared error.
For a noiseless saturated observation \(Y=\{T\geq c\}\), the clipped
estimator is \(a=c\), whereas generally
\(
\mu_c=\mathbb{E}[T\mid T\geq c]>c
\).
Consequently,
\[
\mathbb{E}\!\left[(T-\mu_c)^2\mid T\geq c\right]
\leq
\mathbb{E}\!\left[(T-c)^2\mid T\geq c\right],
\]
with strict inequality unless \(c=\mu_c\). With measurement noise, the same
statement holds after replacing the conditioning event by
\(\{T+\varepsilon\geq c\}\).

The argument also extends through an affine one-step transition. Suppose
\[
T_n=b+B T_{n-1}+\eta_n,
\qquad
\mathbb{E}[\eta_n\mid Y_{n-1}]=0,
\]
where \(B\) is fixed given the model parameters. Then
\[
\mathbb{E}[T_n\mid Y_{n-1}]
=b+B\,\mathbb{E}[T_{n-1}\mid Y_{n-1}].
\]
For a candidate previous-state estimate \(a(Y_{n-1})\), the conditional
squared Euclidean prediction error satisfies
\[
\mathbb{E}\!\left[
\left\|T_n-(b+B a)\right\|^2
\mid Y_{n-1}
\right]
=
\operatorname{tr}\!\left\{\operatorname{Var}(T_n\mid Y_{n-1})\right\}
+
\left\|B\left(a-\mathbb{E}[T_{n-1}\mid Y_{n-1}]\right)\right\|^2.
\]
Thus propagating the previous conditional mean is optimal within this affine
model under squared loss, and improves on clipped propagation whenever the
clipped-state error survives under \(B\). This is a model-based statement: it
requires finite second moments, a correctly specified affine conditional mean,
and zero conditional-mean transition error. It does not guarantee improvement
under transition misspecification or for non-squared probabilistic scores.
